Figure~\ref{fig:fig-auth-path} follows one authorization through six stages:
the account decides at one point, writes nothing on a refusal, and calls the
verifier only after the decision accepts.

\perchfigure{fig-auth-path}

\paragraph{The decision function.} Stage~4 is rule-core's
\code{rule\_eval::check\_auth}~\selfcite{rule-core}
(\nolinkurl{rule-eval/src/authorize.rs:239}):
\begin{Verbatim}[fontsize=\scriptsize,xleftmargin=1em]
pub fn check_auth(rules: &[Rule], records: &[SpendRecord], contexts: &[Context],
    choices: &[u32], authenticated: &[Signer], now: u32, generation: u64)
    -> Result<Vec<SpendRecord>, Refusal>
\end{Verbatim}
It returns the next spending records or a \code{Refusal}, reads no storage and
checks no signature (Definition~\ref{def:check-auth-and-runs} is its model). The
account reports its refusals as auth error \fact{auth-error-stale-grant} (a
stale grant) or \fact{auth-error-rule-eval-refused} (any other), inside
\fact{host-auth-error}, beside the other stages' exits in
Fig.~\ref{fig:fig-auth-path}; a wallet that needs the reason re-runs the
function off chain.

\paragraph{What a signature binds.} An external signer signs a digest of the
host's payload, this account and the selected rule IDs. A delegated signer
(CAP-0071) approves the host's payload, which omits the rule IDs, so a submitter
may select any rule that lists the delegate; OpenZeppelin \fact{oz-v080-tag}
binds them for both kinds.
